% !TEX encoding = UTF-8 Unicode
% !TEX TS-program = pdflatex

% toptesti document class
\documentclass[%
    a4paper, % not needed, by default it is a4paper, or also b5paper can be used
    corpo=12pt, % dimension of basic font
    % oneside is generally the way to go
    oneside, % two side optimizes for two-face printing, having chapters open on the right (aka odd numbers), if you don't want blank pages put oneside here
    stile=standard,
    %evenboxes, % not needed, to put supervisors and candidate at the same level
    tipotesi=magistrale,
    numerazioneromana, % roman numbering for appendixes and preambles, up to Table of Contents
    openright, % to force opening on the right for double-sided printing
    cucitura=7mm, % for printing, 7mm should be enough
    %dvipsnames, % for compatibility with xcolor, it does not work
]{toptesi}

\usepackage{calc} % needed for complex computationsin latex language
% \usepackage[a4paper]{geometry} % needed to change margins


% configuration for glossaries
% convert and load converted glossaries in .tex ,format from .bib
\setabbreviationstyle{long-short-desc} % style before loading resources
% this command sets the style to title for long names of acronyms only in the glossary description, leading to capitalized first-letter for all words
% \glssetcategoryattribute{\glsxtrabbrvtype}{glossname}{capitalisewords} % doesn't work
% resources to load if using a bib file with bib2gls
%\GlsXtrLoadResources[%
% src={glossaries}, % name of the file without extension
% selection=all, % select all the entries
%]
% not needed
%\newglossary*{abbreviation}{Acronyms} % to change the name of this glossary for acronyms

%\renewcommand{\glsclearpage}{\paginavuota} % to allow glossaries to clear pages, done manually is better


% setup for hyperref
\hypersetup{%
    pdfpagemode={UseOutlines},
    bookmarksopen,
    pdfstartview={FitH},
    colorlinks,
    linkcolor={black}, % it is suggested to keep them black, since when printing it it costs per page, and if they have color it's twice the price per page
    citecolor={black},
    urlcolor={black}
  }
%

% setup for svg
\svgsetup{%
    inkscapeformat=pdf, % to force usage of PDF
    inkscapelatex=false, % to disable latex rendering of text, produces errors
}

% setup for siunitx, it does not work in the summary
\sisetup{%
    detect-all, % to use the same font as for writing when using \num
    mode=text, % to allow it to work also in math mode
    group-separator = {,}, % separator for number grouping
    group-minimum-digits = 3, % minimum number of digits a number must have to be grouped in 3-digit groups
}

% listings colours
\definecolor{rulecolor}{rgb}{0,0,0}
\definecolor{commentcolor}{rgb}{0,0.6,0}
\definecolor{linenumbercolor}{rgb}{0.5,0.5,0.5}
\definecolor{keywordcolor}{rgb}{0,0,0.95}
\definecolor{backcolor}{rgb}{1,1,1}%{0.95,0.95,0.92}
\definecolor{stringcolor}{rgb}{0.58,0,0.82}

% setup for lstlisting
\lstset{ %
	backgroundcolor=\color{backcolor},   % choose the background color; you must add \usepackage{color} or \usepackage{xcolor}; should come as last argument
	basicstyle=\footnotesize,        % the size of the fonts that are used for the code
	breakatwhitespace=false,         % sets if automatic breaks should only happen at whitespace
	breaklines=true,                 % sets automatic line breaking
	captionpos=t,                    % sets the caption-position to bottom
	commentstyle=\color{commentcolor},    % comment style
	extendedchars=true,              % lets you use non-ASCII characters; for 8-bits encodings only, does not work with UTF-8
	frame=single,	                   % adds a frame around the code
	keepspaces=true,                 % keeps spaces in text, useful for keeping indentation of code (possibly needs columns=flexible)
	keywordstyle=\color{keywordcolor},       % keyword style
	%language=VHDL,                 % the language of the code
	numbers=left,                    % where to put the line-numbers; possible values are (none, left, right)
	numbersep=5pt,                   % how far the line-numbers are from the code
	numberstyle=\tiny\color{linenumbercolor}, % the style that is used for the line-numbers
	rulecolor=\color{rulecolor},         % if not set, the frame-color may be changed on line-breaks within not-black text (e.g. comments (green here))
	showspaces=false,                % show spaces everywhere adding particular underscores; it overrides 'showstringspaces'
	showstringspaces=false,          % underline spaces within strings only
	showtabs=false,                  % show tabs within strings adding particular underscores
	stepnumber=1,                    % the step between two line-numbers. If it's 1, each line will be numbered
	stringstyle=\color{stringcolor},     % string literal style
	tabsize=4,	                   % sets default tabsize to 2 spaces
	title=\lstname,                   % show the filename of files included with \lstinputlisting; also try caption instead of title
	inputencoding=utf8,
	literate=
	{á}{{\'a}}1 {é}{{\'e}}1 {í}{{\'\i}}1 {ó}{{\'o}}1 {ú}{{\'u}}1
	{Á}{{\'A}}1 {É}{{\'E}}1 {Í}{{\'I}}1 {Ó}{{\'O}}1 {Ú}{{\'U}}1
	{à}{{\`a}}1 {è}{{\`e}}1 {ì}{{\`\i}}1 {ò}{{\`o}}1 {ù}{{\`u}}1
	{À}{{\`A}}1 {È}{{\'E}}1 {Ì}{{\`I}}1 {Ò}{{\`O}}1 {Ù}{{\`U}}1
	{ä}{{\"a}}1 {ë}{{\"e}}1 {ï}{{\"\i}}1 {ö}{{\"o}}1 {ü}{{\"u}}1
	{Ä}{{\"A}}1 {Ë}{{\"E}}1 {Ï}{{\"I}}1 {Ö}{{\"O}}1 {Ü}{{\"U}}1
	{â}{{\^a}}1 {ê}{{\^e}}1 {î}{{\^\i}}1 {ô}{{\^o}}1 {û}{{\^u}}1
	{Â}{{\^A}}1 {Ê}{{\^E}}1 {Î}{{\^I}}1 {Ô}{{\^O}}1 {Û}{{\^U}}1
	{œ}{{\oe}}1 {Œ}{{\OE}}1 {æ}{{\ae}}1 {Æ}{{\AE}}1 {ß}{{\ss}}1
	{ű}{{\H{u}}}1 {Ű}{{\H{U}}}1 {ő}{{\H{o}}}1 {Ő}{{\H{O}}}1
	{ç}{{\c c}}1 {Ç}{{\c C}}1 {ø}{{\o}}1 {å}{{\r a}}1 {Å}{{\r A}}1
	{€}{{\euro}}1 {£}{{\pounds}}1 {«}{{\guillemotleft}}1
	{»}{{\guillemotright}}1 {ñ}{{\~n}}1 {Ñ}{{\~N}}1 {¿}{{?`}}1
}


% biblatex setup
% generally 9000 is ok, if higher than 10000 it's bad
% If you want to break on URL numbers
\setcounter{biburlnumpenalty}{9000}
% If you want to break on URL lower case letters
\setcounter{biburllcpenalty}{9000}
% If you want to break on URL UPPER CASE letters
\setcounter{biburlucpenalty}{9000}



% how to change Contents to Table of Contents
\addto\captionsenglish{% Replace "english" with the language you use
  \renewcommand{\contentsname}%
    {Table of Contents}%
}

% to change the name of Abbreviations to Acronyms
% not needed if use use entry types and define those
% \renewcommand{\abbreviationsname}{Acronyms}

% to allow line comments in algorithms
\algnewcommand{\LineComment}[1]{\State{} \(\triangleright\) #1}

% to declare abs and norm
\DeclarePairedDelimiter\abs{\lvert}{\rvert}%
\DeclarePairedDelimiter\norm{\lVert}{\rVert}%

% Swap the definition of \abs* and \norm*, so that \abs
% and \norm resizes the size of the brackets, and the 
% starred version does not.
\makeatletter
\let\oldabs\abs{}
\def\abs{\@ifstar{\oldabs}{\oldabs*}}
%
\let\oldnorm\norm{}
\def\norm{\@ifstar{\oldnorm}{\oldnorm*}}
\makeatother


% change this configuration with your info
% if you need fewer or more supervisors you have to change \relatore command by adding or removing lines in the table in toptesi_config
\newcommand{\thesistitle}{TITLE}
\newcommand{\thesisuniversitylogo}{images/logo/logoPoliTo_symbol_only} % choose your logo
\newcommand{\thesiscandidatename}{NAME}
\newcommand{\thesiscandidatesurname}{SURNAME}
\newcommand{\thesissupervisoronetitle}{prof.}
\newcommand{\thesissupervisoronename}{NAME}
\newcommand{\thesissupervisoronesurname}{SURNAME}
\newcommand{\thesissupervisortwotitle}{prof.}
\newcommand{\thesissupervisortwoname}{NAME}
\newcommand{\thesissupervisortwosurname}{SURNAME}
\newcommand{\thesissupervisorthreetitle}{prof.}
\newcommand{\thesissupervisorthreename}{NAME}
\newcommand{\thesissupervisorthreesurname}{SURNAME}
\newcommand{\thesisdate}{MONTH YEAR}
\newcommand{\thesiscourse}{COURSE}
\newcommand{\thesisuniversity}{UNIVERSITY}
\newcommand{\thesislevel}{LEVEL} % master or bachelor
\newcommand{\thesiscandidatetext}{Candidate}
\newcommand{\thesissupervisortext}{Supervisors}


% fontsize is {size}{spacing}\family
\newcommand{\institutionfont}{\fontsize{22}{30}\scshape}
\newcommand{\divisionfont}{\fontsize{16}{20}\rmfamily}
\newcommand{\pretitlefont}{\fontsize{16}{16}\rmfamily}
\newcommand{\customtitlefont}{\fontsize{21}{28}\scshape}% {iwona}{bx}{n}}
\newcommand{\fixednamesfont}{\fontsize{14}{20}\mdseries}
\newcommand{\namesfont}{\fontsize{14}{20}\bfseries}
\newcommand{\footfont}{\fontsize{15}{18}\rmfamily}


\addbibresource{bibliography.bib}

% to load the glossaries, not needed if using bib2gls
% for glossary entry
% @entry{bird,
%     name={bird},
%     description = {feathered animal},
%     see={[see also]{duck,goose}}
% }

% if this bib file does not work, try using \input{file.tex}
% where all the \newabbreviation commands have been inserted
% containing all the definitions

% Gls to capitalize first letter
% GLS for full uppercase
% for abbreviations also
% glsxtrshort for abbreviation
% similar for long, full, and capital configurations, add pl at the end for plurals
% glsentryshort, long, plural (referred to shorts) must be used when in section titles
% glslink to allow the link but use a different text (as for href)


% if you want to use also description for the abbreviations/acronyms, you should use bib2gls and define all the entries in a bib file, which is incompatible with Overleaf
\newacronym{AI}{AI}{artificial intelligence}
\makeglossaries{}

\begin{document}

\title{\vspace*{-5mm}\textbf{\thesistitle}\\Summary} % vspace is needed to shift upwards the title
\date{\thesisdate}
\author{\textbf{Candidate}:\\\thesiscandidatename~\thesiscandidatesurname\\
\textbf{Supervisors}:\\\thesissupervisoronetitle~\thesissupervisoronename~\thesissupervisoronesurname\\
\thesissupervisortwotitle~\thesissupervisortwoname~\thesissupervisortwosurname\\
\thesissupervisorthreetitle~\thesissupervisorthreename~\thesissupervisorthreesurname}


\ateneo{\thesisuniversity} % university name
\logosede[5cm]{\thesisuniversitylogo} % logo, square brackets contain the height

\titolo{\thesistitle} % title
%\sottotitolo{Metodo dei satelliti medicei} % subtitle

% place/remove a slash \\ to put the name on the following line or after Master Degree Course
\corsodilaurea{\thesiscourse} % course name


%~251197 % id number is not needed

\candidato{\thesiscandidatename~\textsc{\thesiscandidatesurname}} % candidate

% using tabular we can have more than 1 supervisor under the same column
\relatore{\tabular{@{}l}%
    \xmakefirstuc{\thesissupervisoronetitle}~\thesissupervisoronename~\textsc{\thesissupervisoronesurname}\\[0.4ex]
    \xmakefirstuc{\thesissupervisortwotitle}~\thesissupervisortwoname~\textsc{\thesissupervisortwosurname}\\[0.4ex]
    \xmakefirstuc{\thesissupervisorthreetitle}~\thesissupervisorthreename~\textsc{\thesissupervisorthreesurname}
    \endtabular}
%\terzorelatore{Ciao}

% in this way we have Academic Year without stile=classica, so without lines
%\sedutadilaurea{\textsc{Academic~Year} 2019-2020}% per la laurea magistrale
% for PoliTo there is only month year
\sedutadilaurea{\thesisdate}% per la laurea magistrale
% PhD
%\esamedidottorato{Novembre 1610}
%\ciclodidottorato{XV}

% offset for binding, the smaller the better
%\setbindingcorrection{3mm}


\english% or \italian (default)

\iflanguage{english}{%
	%\retrofrontespizio{This work is subject to the Creative Commons Licence}

	\CorsoDiLaureaIn{\thesislevel's Degree Course in\space}

	\TesiDiLaurea{\thesislevel's Degree Thesis}

	\InName{in}
	\CandidateName{\xmakefirstuc{\thesiscandidatetext}}% or Candidates
	\AdvisorName{\xmakefirstuc{\thesissupervisortext}}% or Supervisor
	%\TutorName{Tutor}
	%\NomeTutoreAziendale{Internship Tutor}

	%\NomePrimoTomo{First volume}
	%\NomeSecondoTomo{Second Volume}
	%\NomeTerzoTomo{Third Volume}
	%\NomeQuartoTomo{Fourth Volume}
}{}


% front page
% frontespizio can be used for the first page print
% while the custom-made frontpage can be used as hard-cover
% use pdfjoin or pdfseparate to extract or put together the pages if needed
%\frontespizio* % without star the logo is on top
\newgeometry{top=4cm,left=3cm,right=3cm,bottom=4cm,heightrounded}
\begin{titlepage}
\centering
%
{\institutionfont{}\textbf{\MakeUppercase{\thesisuniversity}} \par}
%
\vspace{\stretch{2}} % changing this number and the others changes the proportion
%
{\divisionfont{}\textbf{\thesislevel's Degree in \thesiscourse} \par}
%
\vspace{\stretch{3}}
%
\includegraphics[width=50mm]{\thesisuniversitylogo}\\
%
\vspace{\stretch{4}}
%
{\divisionfont{}\textbf{\thesislevel's Degree Thesis} \par}
%
\vspace{\stretch{3}}
%
{\customtitlefont{}\textbf{\thesistitle} \par}
%
\vspace{\stretch{10}}
%
\makebox[\textwidth]{\null\hfill\def\arraystretch{2}% % to change the spacing change this number
\begin{minipage}[t]{.375\textwidth}\raggedright
    \begin{adjustbox}{width={\textwidth},totalheight={\textheight},keepaspectratio} % with adjustbox it adapts to the lengths of the names, remove it if you want the same font dimension
    \begin{tabular}[t]{@{}l@{}}
        \fixednamesfont \textbf{\thesissupervisortext} \\
        \namesfont \xmakefirstuc{\thesissupervisoronetitle}~\thesissupervisoronename~\MakeUppercase{\thesissupervisoronesurname}\\
        \namesfont \xmakefirstuc{\thesissupervisortwotitle}~\thesissupervisortwoname~\MakeUppercase{\thesissupervisortwosurname}\\
        \namesfont \xmakefirstuc{\thesissupervisorthreetitle}~\thesissupervisorthreename~\MakeUppercase{\thesissupervisorthreesurname}
    \end{tabular}
    \end{adjustbox}
\end{minipage}
%
\hfill
%
\begin{minipage}[t]{.375\textwidth}\raggedleft
\begin{adjustbox}{width={\textwidth},totalheight={\textheight},keepaspectratio} % with adjustbox it adapts to the lengths of the names, remove it if you want the same font dimension
\begin{tabular}[t]{@{}l@{}}
    \fixednamesfont \textbf{\thesiscandidatetext} \\
    \namesfont \thesiscandidatename~\MakeUppercase{\thesiscandidatesurname}
\end{tabular}
\end{adjustbox}
\end{minipage}\hfill\null}\\
%
\vspace{\stretch{5}}
%
{\footfont \textbf{\thesisdate} \par}
%
\end{titlepage}

\restoregeometry
 % custom frontpage
%\retrofrontespizio
% insert text for the back of the front page
% if you insert any remove the following \paginavuota
% either a blank page or a back is needed to have double-sided printing
% pay attention to leave the space for the page

%\paginavuota % clears a page

\frontmatter

% abstract if needed
% \begin{abstract}
%     % abstract, choose between abstract and summary
% \end{abstract}

% to create blank pages for openright in frontmatter
% use one of the following two methods
% 1) use the following three lines
%\phantom{0} % needed otherwise cleardoublepage does not clean the page because it sees it empty
%\cleardoublepage
%\thispagestyle{empty} % to have empty page, without numbers
% 2) or
\paginavuota{} % to manually create a blank page

\sommario{}
% !TEX encoding = UTF-8 Unicode
% !TEX TS-program = pdflatex

\documentclass[%
    12pt, % font size
    oneside, % if it should be done for printing on two-sides or not
    a4paper, % format of paper
    notitlepage, % to remove the title page
]{article}

\usepackage{calc} % needed for complex computationsin latex language
% \usepackage[a4paper]{geometry} % needed to change margins


\usepackage{calc} % needed for complex computationsin latex language
% \usepackage[a4paper]{geometry} % needed to change margins


% configuration for glossaries
% convert and load converted glossaries in .tex ,format from .bib
\setabbreviationstyle{long-short-desc} % style before loading resources
% this command sets the style to title for long names of acronyms only in the glossary description, leading to capitalized first-letter for all words
% \glssetcategoryattribute{\glsxtrabbrvtype}{glossname}{capitalisewords} % doesn't work
% resources to load if using a bib file with bib2gls
%\GlsXtrLoadResources[%
% src={glossaries}, % name of the file without extension
% selection=all, % select all the entries
%]
% not needed
%\newglossary*{abbreviation}{Acronyms} % to change the name of this glossary for acronyms

%\renewcommand{\glsclearpage}{\paginavuota} % to allow glossaries to clear pages, done manually is better


% setup for hyperref
\hypersetup{%
    pdfpagemode={UseOutlines},
    bookmarksopen,
    pdfstartview={FitH},
    colorlinks,
    linkcolor={black}, % it is suggested to keep them black, since when printing it it costs per page, and if they have color it's twice the price per page
    citecolor={black},
    urlcolor={black}
  }
%

% setup for svg
\svgsetup{%
    inkscapeformat=pdf, % to force usage of PDF
    inkscapelatex=false, % to disable latex rendering of text, produces errors
}

% setup for siunitx, it does not work in the summary
\sisetup{%
    detect-all, % to use the same font as for writing when using \num
    mode=text, % to allow it to work also in math mode
    group-separator = {,}, % separator for number grouping
    group-minimum-digits = 3, % minimum number of digits a number must have to be grouped in 3-digit groups
}

% listings colours
\definecolor{rulecolor}{rgb}{0,0,0}
\definecolor{commentcolor}{rgb}{0,0.6,0}
\definecolor{linenumbercolor}{rgb}{0.5,0.5,0.5}
\definecolor{keywordcolor}{rgb}{0,0,0.95}
\definecolor{backcolor}{rgb}{1,1,1}%{0.95,0.95,0.92}
\definecolor{stringcolor}{rgb}{0.58,0,0.82}

% setup for lstlisting
\lstset{ %
	backgroundcolor=\color{backcolor},   % choose the background color; you must add \usepackage{color} or \usepackage{xcolor}; should come as last argument
	basicstyle=\footnotesize,        % the size of the fonts that are used for the code
	breakatwhitespace=false,         % sets if automatic breaks should only happen at whitespace
	breaklines=true,                 % sets automatic line breaking
	captionpos=t,                    % sets the caption-position to bottom
	commentstyle=\color{commentcolor},    % comment style
	extendedchars=true,              % lets you use non-ASCII characters; for 8-bits encodings only, does not work with UTF-8
	frame=single,	                   % adds a frame around the code
	keepspaces=true,                 % keeps spaces in text, useful for keeping indentation of code (possibly needs columns=flexible)
	keywordstyle=\color{keywordcolor},       % keyword style
	%language=VHDL,                 % the language of the code
	numbers=left,                    % where to put the line-numbers; possible values are (none, left, right)
	numbersep=5pt,                   % how far the line-numbers are from the code
	numberstyle=\tiny\color{linenumbercolor}, % the style that is used for the line-numbers
	rulecolor=\color{rulecolor},         % if not set, the frame-color may be changed on line-breaks within not-black text (e.g. comments (green here))
	showspaces=false,                % show spaces everywhere adding particular underscores; it overrides 'showstringspaces'
	showstringspaces=false,          % underline spaces within strings only
	showtabs=false,                  % show tabs within strings adding particular underscores
	stepnumber=1,                    % the step between two line-numbers. If it's 1, each line will be numbered
	stringstyle=\color{stringcolor},     % string literal style
	tabsize=4,	                   % sets default tabsize to 2 spaces
	title=\lstname,                   % show the filename of files included with \lstinputlisting; also try caption instead of title
	inputencoding=utf8,
	literate=
	{á}{{\'a}}1 {é}{{\'e}}1 {í}{{\'\i}}1 {ó}{{\'o}}1 {ú}{{\'u}}1
	{Á}{{\'A}}1 {É}{{\'E}}1 {Í}{{\'I}}1 {Ó}{{\'O}}1 {Ú}{{\'U}}1
	{à}{{\`a}}1 {è}{{\`e}}1 {ì}{{\`\i}}1 {ò}{{\`o}}1 {ù}{{\`u}}1
	{À}{{\`A}}1 {È}{{\'E}}1 {Ì}{{\`I}}1 {Ò}{{\`O}}1 {Ù}{{\`U}}1
	{ä}{{\"a}}1 {ë}{{\"e}}1 {ï}{{\"\i}}1 {ö}{{\"o}}1 {ü}{{\"u}}1
	{Ä}{{\"A}}1 {Ë}{{\"E}}1 {Ï}{{\"I}}1 {Ö}{{\"O}}1 {Ü}{{\"U}}1
	{â}{{\^a}}1 {ê}{{\^e}}1 {î}{{\^\i}}1 {ô}{{\^o}}1 {û}{{\^u}}1
	{Â}{{\^A}}1 {Ê}{{\^E}}1 {Î}{{\^I}}1 {Ô}{{\^O}}1 {Û}{{\^U}}1
	{œ}{{\oe}}1 {Œ}{{\OE}}1 {æ}{{\ae}}1 {Æ}{{\AE}}1 {ß}{{\ss}}1
	{ű}{{\H{u}}}1 {Ű}{{\H{U}}}1 {ő}{{\H{o}}}1 {Ő}{{\H{O}}}1
	{ç}{{\c c}}1 {Ç}{{\c C}}1 {ø}{{\o}}1 {å}{{\r a}}1 {Å}{{\r A}}1
	{€}{{\euro}}1 {£}{{\pounds}}1 {«}{{\guillemotleft}}1
	{»}{{\guillemotright}}1 {ñ}{{\~n}}1 {Ñ}{{\~N}}1 {¿}{{?`}}1
}


% biblatex setup
% generally 9000 is ok, if higher than 10000 it's bad
% If you want to break on URL numbers
\setcounter{biburlnumpenalty}{9000}
% If you want to break on URL lower case letters
\setcounter{biburllcpenalty}{9000}
% If you want to break on URL UPPER CASE letters
\setcounter{biburlucpenalty}{9000}


% configuration for glossaries
% convert and load converted glossaries in .tex ,format from .bib
\setabbreviationstyle{long-short-desc} % style before loading resources
% this command sets the style to title for long names of acronyms only in the glossary description, leading to capitalized first-letter for all words
% \glssetcategoryattribute{\glsxtrabbrvtype}{glossname}{capitalisewords} % doesn't work
% resources to load if using a bib file with bib2gls
%\GlsXtrLoadResources[%
% src={glossaries}, % name of the file without extension
% selection=all, % select all the entries
%]
% not needed
%\newglossary*{abbreviation}{Acronyms} % to change the name of this glossary for acronyms

%\renewcommand{\glsclearpage}{\paginavuota} % to allow glossaries to clear pages, done manually is better


% setup for hyperref
\hypersetup{%
    pdfpagemode={UseOutlines},
    bookmarksopen,
    pdfstartview={FitH},
    colorlinks,
    linkcolor={black}, % it is suggested to keep them black, since when printing it it costs per page, and if they have color it's twice the price per page
    citecolor={black},
    urlcolor={black}
  }
%

% setup for svg
\svgsetup{%
    inkscapeformat=pdf, % to force usage of PDF
    inkscapelatex=false, % to disable latex rendering of text, produces errors
}

% setup for siunitx, it does not work in the summary
\sisetup{%
    detect-all, % to use the same font as for writing when using \num
    mode=text, % to allow it to work also in math mode
    group-separator = {,}, % separator for number grouping
    group-minimum-digits = 3, % minimum number of digits a number must have to be grouped in 3-digit groups
}

% listings colours
\definecolor{rulecolor}{rgb}{0,0,0}
\definecolor{commentcolor}{rgb}{0,0.6,0}
\definecolor{linenumbercolor}{rgb}{0.5,0.5,0.5}
\definecolor{keywordcolor}{rgb}{0,0,0.95}
\definecolor{backcolor}{rgb}{1,1,1}%{0.95,0.95,0.92}
\definecolor{stringcolor}{rgb}{0.58,0,0.82}

% setup for lstlisting
\lstset{ %
	backgroundcolor=\color{backcolor},   % choose the background color; you must add \usepackage{color} or \usepackage{xcolor}; should come as last argument
	basicstyle=\footnotesize,        % the size of the fonts that are used for the code
	breakatwhitespace=false,         % sets if automatic breaks should only happen at whitespace
	breaklines=true,                 % sets automatic line breaking
	captionpos=t,                    % sets the caption-position to bottom
	commentstyle=\color{commentcolor},    % comment style
	extendedchars=true,              % lets you use non-ASCII characters; for 8-bits encodings only, does not work with UTF-8
	frame=single,	                   % adds a frame around the code
	keepspaces=true,                 % keeps spaces in text, useful for keeping indentation of code (possibly needs columns=flexible)
	keywordstyle=\color{keywordcolor},       % keyword style
	%language=VHDL,                 % the language of the code
	numbers=left,                    % where to put the line-numbers; possible values are (none, left, right)
	numbersep=5pt,                   % how far the line-numbers are from the code
	numberstyle=\tiny\color{linenumbercolor}, % the style that is used for the line-numbers
	rulecolor=\color{rulecolor},         % if not set, the frame-color may be changed on line-breaks within not-black text (e.g. comments (green here))
	showspaces=false,                % show spaces everywhere adding particular underscores; it overrides 'showstringspaces'
	showstringspaces=false,          % underline spaces within strings only
	showtabs=false,                  % show tabs within strings adding particular underscores
	stepnumber=1,                    % the step between two line-numbers. If it's 1, each line will be numbered
	stringstyle=\color{stringcolor},     % string literal style
	tabsize=4,	                   % sets default tabsize to 2 spaces
	title=\lstname,                   % show the filename of files included with \lstinputlisting; also try caption instead of title
	inputencoding=utf8,
	literate=
	{á}{{\'a}}1 {é}{{\'e}}1 {í}{{\'\i}}1 {ó}{{\'o}}1 {ú}{{\'u}}1
	{Á}{{\'A}}1 {É}{{\'E}}1 {Í}{{\'I}}1 {Ó}{{\'O}}1 {Ú}{{\'U}}1
	{à}{{\`a}}1 {è}{{\`e}}1 {ì}{{\`\i}}1 {ò}{{\`o}}1 {ù}{{\`u}}1
	{À}{{\`A}}1 {È}{{\'E}}1 {Ì}{{\`I}}1 {Ò}{{\`O}}1 {Ù}{{\`U}}1
	{ä}{{\"a}}1 {ë}{{\"e}}1 {ï}{{\"\i}}1 {ö}{{\"o}}1 {ü}{{\"u}}1
	{Ä}{{\"A}}1 {Ë}{{\"E}}1 {Ï}{{\"I}}1 {Ö}{{\"O}}1 {Ü}{{\"U}}1
	{â}{{\^a}}1 {ê}{{\^e}}1 {î}{{\^\i}}1 {ô}{{\^o}}1 {û}{{\^u}}1
	{Â}{{\^A}}1 {Ê}{{\^E}}1 {Î}{{\^I}}1 {Ô}{{\^O}}1 {Û}{{\^U}}1
	{œ}{{\oe}}1 {Œ}{{\OE}}1 {æ}{{\ae}}1 {Æ}{{\AE}}1 {ß}{{\ss}}1
	{ű}{{\H{u}}}1 {Ű}{{\H{U}}}1 {ő}{{\H{o}}}1 {Ő}{{\H{O}}}1
	{ç}{{\c c}}1 {Ç}{{\c C}}1 {ø}{{\o}}1 {å}{{\r a}}1 {Å}{{\r A}}1
	{€}{{\euro}}1 {£}{{\pounds}}1 {«}{{\guillemotleft}}1
	{»}{{\guillemotright}}1 {ñ}{{\~n}}1 {Ñ}{{\~N}}1 {¿}{{?`}}1
}


% biblatex setup
% generally 9000 is ok, if higher than 10000 it's bad
% If you want to break on URL numbers
\setcounter{biburlnumpenalty}{9000}
% If you want to break on URL lower case letters
\setcounter{biburllcpenalty}{9000}
% If you want to break on URL UPPER CASE letters
\setcounter{biburlucpenalty}{9000}



% how to change Contents to Table of Contents
\addto\captionsenglish{% Replace "english" with the language you use
  \renewcommand{\contentsname}%
    {Table of Contents}%
}

% to change the name of Abbreviations to Acronyms
% not needed if use use entry types and define those
% \renewcommand{\abbreviationsname}{Acronyms}

% to allow line comments in algorithms
\algnewcommand{\LineComment}[1]{\State{} \(\triangleright\) #1}

% to declare abs and norm
\DeclarePairedDelimiter\abs{\lvert}{\rvert}%
\DeclarePairedDelimiter\norm{\lVert}{\rVert}%

% Swap the definition of \abs* and \norm*, so that \abs
% and \norm resizes the size of the brackets, and the 
% starred version does not.
\makeatletter
\let\oldabs\abs{}
\def\abs{\@ifstar{\oldabs}{\oldabs*}}
%
\let\oldnorm\norm{}
\def\norm{\@ifstar{\oldnorm}{\oldnorm*}}
\makeatother


% change this configuration with your info
% if you need fewer or more supervisors you have to change \relatore command by adding or removing lines in the table in toptesi_config
\newcommand{\thesistitle}{TITLE}
\newcommand{\thesisuniversitylogo}{images/logo/logoPoliTo_symbol_only} % choose your logo
\newcommand{\thesiscandidatename}{NAME}
\newcommand{\thesiscandidatesurname}{SURNAME}
\newcommand{\thesissupervisoronetitle}{prof.}
\newcommand{\thesissupervisoronename}{NAME}
\newcommand{\thesissupervisoronesurname}{SURNAME}
\newcommand{\thesissupervisortwotitle}{prof.}
\newcommand{\thesissupervisortwoname}{NAME}
\newcommand{\thesissupervisortwosurname}{SURNAME}
\newcommand{\thesissupervisorthreetitle}{prof.}
\newcommand{\thesissupervisorthreename}{NAME}
\newcommand{\thesissupervisorthreesurname}{SURNAME}
\newcommand{\thesisdate}{MONTH YEAR}
\newcommand{\thesiscourse}{COURSE}
\newcommand{\thesisuniversity}{UNIVERSITY}
\newcommand{\thesislevel}{LEVEL} % master or bachelor
\newcommand{\thesiscandidatetext}{Candidate}
\newcommand{\thesissupervisortext}{Supervisors}


\addbibresource{bibliography.bib}

% to load the glossaries, not needed if using bib2gls
% for glossary entry
% @entry{bird,
%     name={bird},
%     description = {feathered animal},
%     see={[see also]{duck,goose}}
% }

% if this bib file does not work, try using \input{file.tex}
% where all the \newabbreviation commands have been inserted
% containing all the definitions

% Gls to capitalize first letter
% GLS for full uppercase
% for abbreviations also
% glsxtrshort for abbreviation
% similar for long, full, and capital configurations, add pl at the end for plurals
% glsentryshort, long, plural (referred to shorts) must be used when in section titles
% glslink to allow the link but use a different text (as for href)


% if you want to use also description for the abbreviations/acronyms, you should use bib2gls and define all the entries in a bib file, which is incompatible with Overleaf
\newacronym{AI}{AI}{artificial intelligence}
\makeglossaries

\begin{document}

\title{\vspace*{-5mm}\textbf{\thesistitle}\\Summary} % vspace is needed to shift upwards the title
\date{\thesisdate}
\author{\textbf{Candidate}:\\\thesiscandidatename~\thesiscandidatesurname\\
\textbf{Supervisors}:\\\thesissupervisoronetitle~\thesissupervisoronename~\thesissupervisoronesurname\\
\thesissupervisortwotitle~\thesissupervisortwoname~\thesissupervisortwosurname\\
\thesissupervisorthreetitle~\thesissupervisorthreename~\thesissupervisorthreesurname}


\title{\vspace*{-5mm}\textbf{\thesistitle}\\Summary} % vspace is needed to shift upwards the title
\date{\thesisdate}
\author{\textbf{Candidate}:\\\thesiscandidatename~\thesiscandidatesurname\\
\textbf{Supervisors}:\\\thesissupervisoronetitle~\thesissupervisoronename~\thesissupervisoronesurname\\
\thesissupervisortwotitle~\thesissupervisortwoname~\thesissupervisortwosurname\\
\thesissupervisorthreetitle~\thesissupervisorthreename~\thesissupervisorthreesurname}


\maketitle

% !TEX encoding = UTF-8 Unicode
% !TEX TS-program = pdflatex

\documentclass[%
    12pt, % font size
    oneside, % if it should be done for printing on two-sides or not
    a4paper, % format of paper
    notitlepage, % to remove the title page
]{article}

\usepackage{calc} % needed for complex computationsin latex language
% \usepackage[a4paper]{geometry} % needed to change margins


\usepackage{calc} % needed for complex computationsin latex language
% \usepackage[a4paper]{geometry} % needed to change margins


% configuration for glossaries
% convert and load converted glossaries in .tex ,format from .bib
\setabbreviationstyle{long-short-desc} % style before loading resources
% this command sets the style to title for long names of acronyms only in the glossary description, leading to capitalized first-letter for all words
% \glssetcategoryattribute{\glsxtrabbrvtype}{glossname}{capitalisewords} % doesn't work
% resources to load if using a bib file with bib2gls
%\GlsXtrLoadResources[%
% src={glossaries}, % name of the file without extension
% selection=all, % select all the entries
%]
% not needed
%\newglossary*{abbreviation}{Acronyms} % to change the name of this glossary for acronyms

%\renewcommand{\glsclearpage}{\paginavuota} % to allow glossaries to clear pages, done manually is better


% setup for hyperref
\hypersetup{%
    pdfpagemode={UseOutlines},
    bookmarksopen,
    pdfstartview={FitH},
    colorlinks,
    linkcolor={black}, % it is suggested to keep them black, since when printing it it costs per page, and if they have color it's twice the price per page
    citecolor={black},
    urlcolor={black}
  }
%

% setup for svg
\svgsetup{%
    inkscapeformat=pdf, % to force usage of PDF
    inkscapelatex=false, % to disable latex rendering of text, produces errors
}

% setup for siunitx, it does not work in the summary
\sisetup{%
    detect-all, % to use the same font as for writing when using \num
    mode=text, % to allow it to work also in math mode
    group-separator = {,}, % separator for number grouping
    group-minimum-digits = 3, % minimum number of digits a number must have to be grouped in 3-digit groups
}

% listings colours
\definecolor{rulecolor}{rgb}{0,0,0}
\definecolor{commentcolor}{rgb}{0,0.6,0}
\definecolor{linenumbercolor}{rgb}{0.5,0.5,0.5}
\definecolor{keywordcolor}{rgb}{0,0,0.95}
\definecolor{backcolor}{rgb}{1,1,1}%{0.95,0.95,0.92}
\definecolor{stringcolor}{rgb}{0.58,0,0.82}

% setup for lstlisting
\lstset{ %
	backgroundcolor=\color{backcolor},   % choose the background color; you must add \usepackage{color} or \usepackage{xcolor}; should come as last argument
	basicstyle=\footnotesize,        % the size of the fonts that are used for the code
	breakatwhitespace=false,         % sets if automatic breaks should only happen at whitespace
	breaklines=true,                 % sets automatic line breaking
	captionpos=t,                    % sets the caption-position to bottom
	commentstyle=\color{commentcolor},    % comment style
	extendedchars=true,              % lets you use non-ASCII characters; for 8-bits encodings only, does not work with UTF-8
	frame=single,	                   % adds a frame around the code
	keepspaces=true,                 % keeps spaces in text, useful for keeping indentation of code (possibly needs columns=flexible)
	keywordstyle=\color{keywordcolor},       % keyword style
	%language=VHDL,                 % the language of the code
	numbers=left,                    % where to put the line-numbers; possible values are (none, left, right)
	numbersep=5pt,                   % how far the line-numbers are from the code
	numberstyle=\tiny\color{linenumbercolor}, % the style that is used for the line-numbers
	rulecolor=\color{rulecolor},         % if not set, the frame-color may be changed on line-breaks within not-black text (e.g. comments (green here))
	showspaces=false,                % show spaces everywhere adding particular underscores; it overrides 'showstringspaces'
	showstringspaces=false,          % underline spaces within strings only
	showtabs=false,                  % show tabs within strings adding particular underscores
	stepnumber=1,                    % the step between two line-numbers. If it's 1, each line will be numbered
	stringstyle=\color{stringcolor},     % string literal style
	tabsize=4,	                   % sets default tabsize to 2 spaces
	title=\lstname,                   % show the filename of files included with \lstinputlisting; also try caption instead of title
	inputencoding=utf8,
	literate=
	{á}{{\'a}}1 {é}{{\'e}}1 {í}{{\'\i}}1 {ó}{{\'o}}1 {ú}{{\'u}}1
	{Á}{{\'A}}1 {É}{{\'E}}1 {Í}{{\'I}}1 {Ó}{{\'O}}1 {Ú}{{\'U}}1
	{à}{{\`a}}1 {è}{{\`e}}1 {ì}{{\`\i}}1 {ò}{{\`o}}1 {ù}{{\`u}}1
	{À}{{\`A}}1 {È}{{\'E}}1 {Ì}{{\`I}}1 {Ò}{{\`O}}1 {Ù}{{\`U}}1
	{ä}{{\"a}}1 {ë}{{\"e}}1 {ï}{{\"\i}}1 {ö}{{\"o}}1 {ü}{{\"u}}1
	{Ä}{{\"A}}1 {Ë}{{\"E}}1 {Ï}{{\"I}}1 {Ö}{{\"O}}1 {Ü}{{\"U}}1
	{â}{{\^a}}1 {ê}{{\^e}}1 {î}{{\^\i}}1 {ô}{{\^o}}1 {û}{{\^u}}1
	{Â}{{\^A}}1 {Ê}{{\^E}}1 {Î}{{\^I}}1 {Ô}{{\^O}}1 {Û}{{\^U}}1
	{œ}{{\oe}}1 {Œ}{{\OE}}1 {æ}{{\ae}}1 {Æ}{{\AE}}1 {ß}{{\ss}}1
	{ű}{{\H{u}}}1 {Ű}{{\H{U}}}1 {ő}{{\H{o}}}1 {Ő}{{\H{O}}}1
	{ç}{{\c c}}1 {Ç}{{\c C}}1 {ø}{{\o}}1 {å}{{\r a}}1 {Å}{{\r A}}1
	{€}{{\euro}}1 {£}{{\pounds}}1 {«}{{\guillemotleft}}1
	{»}{{\guillemotright}}1 {ñ}{{\~n}}1 {Ñ}{{\~N}}1 {¿}{{?`}}1
}


% biblatex setup
% generally 9000 is ok, if higher than 10000 it's bad
% If you want to break on URL numbers
\setcounter{biburlnumpenalty}{9000}
% If you want to break on URL lower case letters
\setcounter{biburllcpenalty}{9000}
% If you want to break on URL UPPER CASE letters
\setcounter{biburlucpenalty}{9000}


% configuration for glossaries
% convert and load converted glossaries in .tex ,format from .bib
\setabbreviationstyle{long-short-desc} % style before loading resources
% this command sets the style to title for long names of acronyms only in the glossary description, leading to capitalized first-letter for all words
% \glssetcategoryattribute{\glsxtrabbrvtype}{glossname}{capitalisewords} % doesn't work
% resources to load if using a bib file with bib2gls
%\GlsXtrLoadResources[%
% src={glossaries}, % name of the file without extension
% selection=all, % select all the entries
%]
% not needed
%\newglossary*{abbreviation}{Acronyms} % to change the name of this glossary for acronyms

%\renewcommand{\glsclearpage}{\paginavuota} % to allow glossaries to clear pages, done manually is better


% setup for hyperref
\hypersetup{%
    pdfpagemode={UseOutlines},
    bookmarksopen,
    pdfstartview={FitH},
    colorlinks,
    linkcolor={black}, % it is suggested to keep them black, since when printing it it costs per page, and if they have color it's twice the price per page
    citecolor={black},
    urlcolor={black}
  }
%

% setup for svg
\svgsetup{%
    inkscapeformat=pdf, % to force usage of PDF
    inkscapelatex=false, % to disable latex rendering of text, produces errors
}

% setup for siunitx, it does not work in the summary
\sisetup{%
    detect-all, % to use the same font as for writing when using \num
    mode=text, % to allow it to work also in math mode
    group-separator = {,}, % separator for number grouping
    group-minimum-digits = 3, % minimum number of digits a number must have to be grouped in 3-digit groups
}

% listings colours
\definecolor{rulecolor}{rgb}{0,0,0}
\definecolor{commentcolor}{rgb}{0,0.6,0}
\definecolor{linenumbercolor}{rgb}{0.5,0.5,0.5}
\definecolor{keywordcolor}{rgb}{0,0,0.95}
\definecolor{backcolor}{rgb}{1,1,1}%{0.95,0.95,0.92}
\definecolor{stringcolor}{rgb}{0.58,0,0.82}

% setup for lstlisting
\lstset{ %
	backgroundcolor=\color{backcolor},   % choose the background color; you must add \usepackage{color} or \usepackage{xcolor}; should come as last argument
	basicstyle=\footnotesize,        % the size of the fonts that are used for the code
	breakatwhitespace=false,         % sets if automatic breaks should only happen at whitespace
	breaklines=true,                 % sets automatic line breaking
	captionpos=t,                    % sets the caption-position to bottom
	commentstyle=\color{commentcolor},    % comment style
	extendedchars=true,              % lets you use non-ASCII characters; for 8-bits encodings only, does not work with UTF-8
	frame=single,	                   % adds a frame around the code
	keepspaces=true,                 % keeps spaces in text, useful for keeping indentation of code (possibly needs columns=flexible)
	keywordstyle=\color{keywordcolor},       % keyword style
	%language=VHDL,                 % the language of the code
	numbers=left,                    % where to put the line-numbers; possible values are (none, left, right)
	numbersep=5pt,                   % how far the line-numbers are from the code
	numberstyle=\tiny\color{linenumbercolor}, % the style that is used for the line-numbers
	rulecolor=\color{rulecolor},         % if not set, the frame-color may be changed on line-breaks within not-black text (e.g. comments (green here))
	showspaces=false,                % show spaces everywhere adding particular underscores; it overrides 'showstringspaces'
	showstringspaces=false,          % underline spaces within strings only
	showtabs=false,                  % show tabs within strings adding particular underscores
	stepnumber=1,                    % the step between two line-numbers. If it's 1, each line will be numbered
	stringstyle=\color{stringcolor},     % string literal style
	tabsize=4,	                   % sets default tabsize to 2 spaces
	title=\lstname,                   % show the filename of files included with \lstinputlisting; also try caption instead of title
	inputencoding=utf8,
	literate=
	{á}{{\'a}}1 {é}{{\'e}}1 {í}{{\'\i}}1 {ó}{{\'o}}1 {ú}{{\'u}}1
	{Á}{{\'A}}1 {É}{{\'E}}1 {Í}{{\'I}}1 {Ó}{{\'O}}1 {Ú}{{\'U}}1
	{à}{{\`a}}1 {è}{{\`e}}1 {ì}{{\`\i}}1 {ò}{{\`o}}1 {ù}{{\`u}}1
	{À}{{\`A}}1 {È}{{\'E}}1 {Ì}{{\`I}}1 {Ò}{{\`O}}1 {Ù}{{\`U}}1
	{ä}{{\"a}}1 {ë}{{\"e}}1 {ï}{{\"\i}}1 {ö}{{\"o}}1 {ü}{{\"u}}1
	{Ä}{{\"A}}1 {Ë}{{\"E}}1 {Ï}{{\"I}}1 {Ö}{{\"O}}1 {Ü}{{\"U}}1
	{â}{{\^a}}1 {ê}{{\^e}}1 {î}{{\^\i}}1 {ô}{{\^o}}1 {û}{{\^u}}1
	{Â}{{\^A}}1 {Ê}{{\^E}}1 {Î}{{\^I}}1 {Ô}{{\^O}}1 {Û}{{\^U}}1
	{œ}{{\oe}}1 {Œ}{{\OE}}1 {æ}{{\ae}}1 {Æ}{{\AE}}1 {ß}{{\ss}}1
	{ű}{{\H{u}}}1 {Ű}{{\H{U}}}1 {ő}{{\H{o}}}1 {Ő}{{\H{O}}}1
	{ç}{{\c c}}1 {Ç}{{\c C}}1 {ø}{{\o}}1 {å}{{\r a}}1 {Å}{{\r A}}1
	{€}{{\euro}}1 {£}{{\pounds}}1 {«}{{\guillemotleft}}1
	{»}{{\guillemotright}}1 {ñ}{{\~n}}1 {Ñ}{{\~N}}1 {¿}{{?`}}1
}


% biblatex setup
% generally 9000 is ok, if higher than 10000 it's bad
% If you want to break on URL numbers
\setcounter{biburlnumpenalty}{9000}
% If you want to break on URL lower case letters
\setcounter{biburllcpenalty}{9000}
% If you want to break on URL UPPER CASE letters
\setcounter{biburlucpenalty}{9000}



% how to change Contents to Table of Contents
\addto\captionsenglish{% Replace "english" with the language you use
  \renewcommand{\contentsname}%
    {Table of Contents}%
}

% to change the name of Abbreviations to Acronyms
% not needed if use use entry types and define those
% \renewcommand{\abbreviationsname}{Acronyms}

% to allow line comments in algorithms
\algnewcommand{\LineComment}[1]{\State{} \(\triangleright\) #1}

% to declare abs and norm
\DeclarePairedDelimiter\abs{\lvert}{\rvert}%
\DeclarePairedDelimiter\norm{\lVert}{\rVert}%

% Swap the definition of \abs* and \norm*, so that \abs
% and \norm resizes the size of the brackets, and the 
% starred version does not.
\makeatletter
\let\oldabs\abs{}
\def\abs{\@ifstar{\oldabs}{\oldabs*}}
%
\let\oldnorm\norm{}
\def\norm{\@ifstar{\oldnorm}{\oldnorm*}}
\makeatother


% change this configuration with your info
% if you need fewer or more supervisors you have to change \relatore command by adding or removing lines in the table in toptesi_config
\newcommand{\thesistitle}{TITLE}
\newcommand{\thesisuniversitylogo}{images/logo/logoPoliTo_symbol_only} % choose your logo
\newcommand{\thesiscandidatename}{NAME}
\newcommand{\thesiscandidatesurname}{SURNAME}
\newcommand{\thesissupervisoronetitle}{prof.}
\newcommand{\thesissupervisoronename}{NAME}
\newcommand{\thesissupervisoronesurname}{SURNAME}
\newcommand{\thesissupervisortwotitle}{prof.}
\newcommand{\thesissupervisortwoname}{NAME}
\newcommand{\thesissupervisortwosurname}{SURNAME}
\newcommand{\thesissupervisorthreetitle}{prof.}
\newcommand{\thesissupervisorthreename}{NAME}
\newcommand{\thesissupervisorthreesurname}{SURNAME}
\newcommand{\thesisdate}{MONTH YEAR}
\newcommand{\thesiscourse}{COURSE}
\newcommand{\thesisuniversity}{UNIVERSITY}
\newcommand{\thesislevel}{LEVEL} % master or bachelor
\newcommand{\thesiscandidatetext}{Candidate}
\newcommand{\thesissupervisortext}{Supervisors}


\addbibresource{bibliography.bib}

% to load the glossaries, not needed if using bib2gls
% for glossary entry
% @entry{bird,
%     name={bird},
%     description = {feathered animal},
%     see={[see also]{duck,goose}}
% }

% if this bib file does not work, try using \input{file.tex}
% where all the \newabbreviation commands have been inserted
% containing all the definitions

% Gls to capitalize first letter
% GLS for full uppercase
% for abbreviations also
% glsxtrshort for abbreviation
% similar for long, full, and capital configurations, add pl at the end for plurals
% glsentryshort, long, plural (referred to shorts) must be used when in section titles
% glslink to allow the link but use a different text (as for href)


% if you want to use also description for the abbreviations/acronyms, you should use bib2gls and define all the entries in a bib file, which is incompatible with Overleaf
\newacronym{AI}{AI}{artificial intelligence}
\makeglossaries

\begin{document}

\title{\vspace*{-5mm}\textbf{\thesistitle}\\Summary} % vspace is needed to shift upwards the title
\date{\thesisdate}
\author{\textbf{Candidate}:\\\thesiscandidatename~\thesiscandidatesurname\\
\textbf{Supervisors}:\\\thesissupervisoronetitle~\thesissupervisoronename~\thesissupervisoronesurname\\
\thesissupervisortwotitle~\thesissupervisortwoname~\thesissupervisortwosurname\\
\thesissupervisorthreetitle~\thesissupervisorthreename~\thesissupervisorthreesurname}


\title{\vspace*{-5mm}\textbf{\thesistitle}\\Summary} % vspace is needed to shift upwards the title
\date{\thesisdate}
\author{\textbf{Candidate}:\\\thesiscandidatename~\thesiscandidatesurname\\
\textbf{Supervisors}:\\\thesissupervisoronetitle~\thesissupervisoronename~\thesissupervisoronesurname\\
\thesissupervisortwotitle~\thesissupervisortwoname~\thesissupervisortwosurname\\
\thesissupervisorthreetitle~\thesissupervisorthreename~\thesissupervisorthreesurname}


\maketitle

% !TEX encoding = UTF-8 Unicode
% !TEX TS-program = pdflatex

\documentclass[%
    12pt, % font size
    oneside, % if it should be done for printing on two-sides or not
    a4paper, % format of paper
    notitlepage, % to remove the title page
]{article}

\input{common/packages}

\input{common/summary/packages.tex}

\input{common/package_config}

\input{common/summary/package_config.tex}

\input{common/new_commands}

\input{common/thesis_info.tex}

\addbibresource{bibliography.bib}

% to load the glossaries, not needed if using bib2gls
\input{glossaries.tex}
\makeglossaries

\begin{document}

\input{common/config}

\input{common/summary/config.tex}

\maketitle

\input{content/summary.tex}

\end{document}


\end{document}


\end{document}


\phantom{0}
\cleardoublepage{}
\thispagestyle{empty}

\ringraziamenti{} % acknowledgements
% acknowledgements

ACKNOWLEDGMENTS

\vspace*{5\baselineskip}

\begin{flushright}
    \textit{``HI''\\
    Goofy, \href{https://google.com}{Google} by Google}
\end{flushright}


\paginavuota{}
\tableofcontents

\listoftables % ToC for tables

\listoffigures % ToC for figures

% actually abbreviation is the name used for acronym in glossaries-extra
% title sets the name
% type tells the type of glossary to print
% style overrides the global style
% here we are printing only abbreviations
% printunsrtglossary if using record, otherwise printglossary is ok
\paginavuota{}
\printunsrtglossary[style=altlist,title=Acronyms,type=\glsxtrabbrvtype]

% also list of symbols here if needed

% to remove all first use occurrences given the presence of the summary
\glsresetall
% to skip all the first use occurrences, using only short forms
% \glsunsetall


\mainmatter{}

%\part{Prima Parte} % parts division, not needed
% Chapters always open on a right-side page, i.e. odd numbers, so a blank page is inserted if needed
%\cleardoublepage[empty] % to have a fully blank page
% a blank page appears before the first chapter in some configurations, on the last version it doesn't

% list here all the chapters
\chapter{Hello}
\label{sec:hello}

% use [][] to prepend/postpone text to the citation
\cite[Hi][Goofy]{IEEEexample:article_typical}

\si{\kilo\gram\per\second}

% generic figure
\begin{figure}[h]
\centering
\includegraphics[width=.9\linewidth]{images/logo/logoPoliTo_with_name_wrong.png}
\caption{Hi}
\label{fig:hi}
\end{figure}

% use [] to set name for ToC
\section[Extremely long name with manual linebreak which otherwise would not fit the page]{Extremely long name with manual linebreak\\which otherwise would not fit the page} % ok with fontsize=12pt

% list
\begin{enumerate}
    \item A
    \item B
    \item C
\end{enumerate}

% minipage to put two images in the same figure
\begin{figure}[h]
    \centering
    \begin{minipage}[t]{.49\linewidth}
    \begin{figure}[H]
	\centering
	\includegraphics[width=\linewidth]{images/logo/logoPoliTo_with_name_low_quality.jpg}
	\caption{HI}
	\label{fig:c}
    \end{figure}
    \end{minipage}
    \hfill
    \begin{minipage}[t]{.49\linewidth}
    \begin{figure}[H]
	\centering
	% svg inclusion, requires inkscape
	\includesvg[width=\linewidth]{images/artificial_neural_network.svg}
	\caption{SVG}
	\label{fig:svg}
    \end{figure}
    \end{minipage}
\end{figure}

\begin{table}[]
    \centering
    \setcellgapes{3pt}
    \makegapedcells
    \begin{tabular}{|c|c|c}
    \hline
    ReLU & $f(x) = \begin{cases}
	0 & \text{for } x \le 0\\
	x & \text{for } x > 0\end{cases}$ \\ \hline
    Softmax & $f_i(\vec{x}) = \dfrac{e^{x_i}}{\sum_{j=1}^J e^{x_j}} i = 1, ..., J$ \\ \hline
    tanh & $f(x)=\tanh(x)=\dfrac{(e^{x} - e^{-x})}{(e^{x} + e^{-x})}$ \\ \hline
    \end{tabular}
    \caption{Examples of activation functions, operating either element-wise or vector-wise, depending on the function}
    \label{tab:activation_functions}
\end{table}

\begin{equation}
    \label{eq:fully_connected}
    output = f_{activation}\left(\displaystyle\sum_{\#neurons} input_i + bias\right)
\end{equation}

\begin{table}
    \centering
    \begin{adjustbox}{width={0.9\textwidth},totalheight={\textheight},keepaspectratio} % needed if the table overflows the margins, requires adjustbox package
    \setcellgapes{3pt}
    \makegapedcells
    \begin{tabular}{|c|c|}
    \hline
    MSE / L2 Loss / Quadratic Loss & $\dfrac{\sum_{i=1}^{N} \left(y_i - \hat{y}_i\right)^2}{N}$ \\ \hline
    \makecell{(Binary) Cross Entropy \\ (average reduction on higher dimensions)} & $\dfrac{\sum_{i=1}^{N} \sum_{j=1}^{C} \hat{y}_i \log\left(y_{i,j}\right)}{N}$ \\ \hline
    \makecell{Categorical Cross Entropy \\ (sum reduction on higher dimensions)} & $- \sum_{i=1}^{N} \hat{y}_i +  \log\left(\sum_{i=1}^{N} \sum_{j=1}^{C} y_{i,j}\right)$ \\ \hline
    \end{tabular}
    \end{adjustbox} % must be closed before label and caption
    \caption{$y$ is the output of the network, $N$ is the batch size multiplied by the number of outputs (e.g. pixels), $C$ is the number of classes and $\hat{y}$ is the correct output.}
    \label{tab:loss_functions}
\end{table}


\begin{algorithm}
\caption{Adam optimizer algorithm. All operations are element-wise, even powers. Good values for the constants are $\alpha=0.001, \beta_1 = 0.9, \beta_2 = 0.999, \epsilon = 10^{-8}$. $\epsilon$ is needed to guarantee numerical stability.}
\label{alg:adam_optimizer}
\begin{algorithmic}[1]
\Procedure{Adam}{$\alpha, \beta_1, \beta_2, f, \theta_0$}
\LineComment{$\alpha$ is the stepsize}
\LineComment{$\beta_1, \beta_2 \in \left[0, 1\right)$ are the exponential decay rates for the moment estimates}
\LineComment{$f\left(\theta\right)$ is the objective function to optimize}
\LineComment{$\theta_0$ is the initial vector of parameters which will be optimized}
\LineComment{Initialization}
\State $m_0 \gets 0$
\Comment{First moment estimate vector set to 0}
\State $v_0 \gets 0$
\Comment{Second moment estimate vector set to 0}
\State $t \gets 0$
\Comment{Timestep set to 0}
\LineComment{Execution}
\While{$\theta_t$ not converged}
\State $t \gets t + 1$
\Comment{Update timestep}
\LineComment{Gradients are computed w.r.t the parameters to optimize}
\LineComment{using the value of the objective function}
\LineComment{at the previous timestep}
\State $g_t \gets \nabla_\theta f\left(\theta_{t - 1}\right)$
\LineComment{Update of first-moment and second-moment estimates using}
\LineComment{previous value and new gradients, biased}
\State $m_t \gets \beta_1 \cdot m_{t - 1} + \left( 1 - \beta_1 \right) \cdot g_t$
\State $v_t \gets \beta_2 \cdot v_{t - 1} + \left(1 - \beta_2 \right) \cdot g_t^2$
\LineComment{Bias-correction of estimates}
\State $\hat{m}_t \gets \dfrac{m_t}{1 - \beta_1^t}$
\State $\hat{v}_t \gets \dfrac{v_t}{1 - \beta_2^t}$
\State $\theta_t \gets \theta_{t - 1} - \alpha \cdot \dfrac{\hat{m}_t}{\sqrt{\hat{v}_t} + \epsilon}$
\Comment{Update parameters}
\EndWhile
\State \textbf{return} $\theta_t$
\Comment{Optimized parameters are returned}
\EndProcedure
%\end{small}
\end{algorithmic}
\end{algorithm}

% bullet points
\begin{itemize}
    \item A
    \item B
    \item C
\end{itemize}



% \paginavuota % it works even without stile=classica

\appendix
% appendix
\chapter{Galileo}
\label{sec:appendix_galileo}

%\lstinputlisting[]{} % for source code files directly
% lstlisting environment for direct inclusion
\begin{lstlisting}[language=Python]
    import os
    os.system("echo 1")
\end{lstlisting}

% for computational complexity
$\mathcal{O}\left(n\log{n}\right)$

% verbatim
\verb+numpy+



% endnotes here if needed

\phantom{0}
\cleardoublepage{}
\printbibliography[heading=bibintoc] % heading required to show it in ToC

\end{document}
